\documentclass[a4paper]{article}

\usepackage[spanish]{babel}
\usepackage{graphicx}
\usepackage{amsmath, amssymb}
\usepackage[margin=2cm]{geometry}
\usepackage{fancyhdr}
\usepackage{enumerate}
\usepackage[shortlabels]{enumitem}
\usepackage{parskip}
\usepackage[most]{tcolorbox}
\usepackage[hidelinks]{hyperref}
\usepackage{float}
\usepackage{booktabs}



% cabecera
\pagestyle{fancy}
\fancyhead[l]{Mecánica Celeste}
\fancyhead[c]{Clase \#1}
\fancyhead[r]{\today}
\fancyfoot[c]{\thepage}
\renewcommand{\headrulewidth}{0.2pt} % linea horizontal


\begin{document}

\section{Presentación del curso}
La \textbf{Mecánica Celeste} es la disciplina científica encargada de aplicar las leyes de la mecánica para estudiar el movimiento de cuerpos bajo la acción dominante de la \textbf{gravedad}. Aunque clásicamente se definía como el estudio de las posiciones de los planetas, hoy abarca desde pequeñas partículas de polvo hasta galaxias y vehículos espaciales. Este curso está dirigido a estudiantes de pregrado en astronomía, física e ingeniería aeroespacial.

\subsection{Raíces históricas}
El descubrimiento de los planetas se remonta a la antigüedad, cuando se observaron cuerpos que no mantenían posiciones fijas entre las estrellas. Civilizaciones como los babilonios y asirios fueron pioneras en predecir posiciones planetarias usando métodos aritméticos y geométricos. Durante siglos, modelos como las esferas de Eudoxio y el sistema de Claudio Ptolomeo dominaron la disciplina.

\section{Estructura y Contenido}
El curso se organiza en tres grandes bloques temáticos que siguen un orden lógico y atemporal:
\begin{enumerate}
    \item \textbf{Fundamentos:} Repaso de bases matemáticas (cálculo vectorial, cónicas) y físicas (mecánica newtoniana).
    \item \textbf{Formalismo Vectorial:} Estudio del problema de los $N$-cuerpos y el problema relativo de dos cuerpos.
    \item \textbf{Mecánica Analítica:} Introducción a los formalismos Lagrangiano y Hamiltoniano aplicados a la dinámica celeste.
\end{enumerate}

\section{Recursos y Metodología}
El curso integra herramientas modernas para el aprendizaje y la investigación:
\begin{itemize}
    \item \textbf{Texto guía:} \textit{Mecánica Celeste: teoría, algoritmos y problemas}, del Prof. Jorge I. Zuluaga.
    \item \textbf{Lenguaje de programación:} Se utilizará \textbf{Python} por su sintaxis cercana al lenguaje natural y su amplia comunidad científica.
    \item \textbf{PymCel:} Paquete de utilidades diseñado originalmente para el libro guía que facilita el estudio de la astrodinámica.
    \item \textbf{SPICE:} Sistema de la NASA para la determinación precisa de posiciones de cuerpos del Sistema Solar y naves espaciales.
    \item \textbf{Jupyter Notebooks:} El material está disponible en libretas interactivas que permiten ejecutar algoritmos y visualizaciones en clase.
\end{itemize}

\section{Evaluación}
El sistema de evaluación se basa en el desarrollo de cuatro competencias principales:
\begin{itemize}
    \item \textbf{Asistencia:} Seguimiento con nota (10\%).
    \item \textbf{Conceptos:} Evaluados mediante quices (25\%) y un examen final (20\%).
    \item \textbf{Modelación teórica:} Solución de problemas teóricos (30\%).
    \item \textbf{Modelación computacional:} Desarrollo de un proyecto computacional (15\%).
\end{itemize}

\section{El Problema Básico de la Mecánica Celeste}
El objetivo fundamental es predecir la posición y velocidad de los cuerpos celestes en el futuro dadas sus condiciones iniciales. Para un sistema de $N$ partículas puntuales que interactúan gravitacionalmente, el movimiento se describe mediante un conjunto de ecuaciones diferenciales ordinarias.

\subsection{Ecuaciones de movimiento}
Para cada partícula $i$ en un sistema, la aceleración se determina por la suma de las fuerzas gravitacionales ejercidas por las demás partículas $j$:
\begin{equation}
    \left\{ \ddot{\vec{r}}_i = -\sum_{j=1, j \neq i}^{N} \frac{\mu_j}{r_{ij}^3} \vec{r}_{ij} \right\}_{i=1...N}
\end{equation}
Donde $\mu_j = Gm_j$ es el parámetro gravitacional del cuerpo $j$ y $\vec{r}_{ij} = \vec{r}_i - \vec{r}_j$ representa el vector de posición relativa entre los cuerpos. En un sistema de dos cuerpos como la Tierra y el Sol, ambos experimentan fuerzas de atracción mutua, lo que implica que el Sol también está en movimiento y el sistema debe analizarse de forma relativa o respecto al centro de masa.


\bibliographystyle{plainurl}
\bibliography{bibliografia} 
\end{document}
